\chapter{Introduction}\label{introduction}


The introduction describes the problem, why it is important, the main
ideas of the thesis, what are the main contributions of the
thesis, etc. The introduction is the most important part of your thesis because it answers the question of why you are researching your topic at all. The introduction does not only motivate your thesis, it also justifies it. 
Compared to the rest of your thesis, it is a very short text, but this text should be extremely well-written and clear. It deserves a lot of polishing. If the introduction is boring or not well written, a reader (and your supervisor who will finally grade your thesis) will most likely not have fun reading your work. The introduction (together with abstract and title) is the ``bait'' to get your reader, your supervisor and your reviewer hooked.

The introduction should be a short version of the whole thesis, with a focus on motivation. A good introduction is structured as follows:

\begin{enumerate}
\item It should start with a brief motivational problem outline to set the stage, e.g., ``It is important to build intelligent robots. A popular method to build intelligent robots is X. However, X has problem Y.''

\item It should then contain a specific formulation of your research question (RQ). For a good description of how to formulate a research question, consider \url{https://www.scribbr.com/research-process/research-questions/}. E.g.: ``Problem Y has not yet been addressed sufficiently by other researchers, so how can we extend X to address Y?''

\item In addition to your research question, you can (but don't have to) state a specific hypothesis: 
``I hypothesize that we can address Y by extending X with Z.'' It is important that you also say how you plan to validate your hypothesis. For example, if you do an empirical investigation, write something like ``I will validate my hypothesis by measuring ... with experiment ...''. If you write a more formal thesis, you might want to prove a certain theorem, so you could write ``I will validate my hypothesis by proving that ...''

\item The RQ and hypothesis should be followed by a very brief (one-two sentences) summary of how you want to address the question, i.e., the methods and evaluation metrics you want to use, highlighting your research contributions with respect to the state of the art. 
Herein, you should also refer to the individual sections in your paper to give the reader a brief outline. 
``The review of the related work in Section 2 of this paper shows that method Q has not yet been investigated as a solution to problem X. As a novel contribution, I will address this hypothesis / RQ by implementing Z using method Q, as described in Section 3. I will evaluate the results by measuring V, as described in Section 4.'' (Note that you might not have a Section 4 if you do not conduct empirical research  but focus on theory instead.)

\item You should conclude with a sentence like this: ``I conclude in Section 5, by highlighting that extending X with Z indeed achieves Y, but only if W holds.'' 

\end{enumerate}
It does not matter whether you use ``I'' or ``We''. It does matter, however, that you use the same pronoun consistently.

